% !TEX root = ../main.tex

\begin{digest}

Researchers use search engines, reference managers, language-model assistants, and citation graphs throughout a literature workflow. These tools help at different stages, yet a reader still has to preserve the relationship among a changing paper, a generated explanation, and the personal context in which the paper became relevant. MNEME addresses this problem through a mobile-native literature loop for receiving papers, reading evidence-linked explanations, recording bounded interest signals, preparing a research briefing, and exploring related work. The system treats paper identity, evidence, behavior, and mobile state as explicit contracts shared by the Android client and the backend.

The design is grounded in ten semi-structured interviews with students and early-stage researchers. Within this formative sample, eight participants described difficulty recovering a previously encountered paper, seven described paper-volume overload, seven described weak support for connecting papers with developing ideas, and five raised concerns about AI-generated material without inspectable sources. Affinity mapping and customer-profile synthesis translated these observations into requirements for bounded briefings, stable paper identity, visible source actions, correctable personalization, a limited citation neighborhood, and explicit live, cached, processing, fallback, and error states. A five-participant prototype study found full completion of briefing, source, and cited-question tasks, four completions for interest editing, and two for graph exploration. The graph finding led to clearer instructions, persistent node selection, relation context, and a visible \emph{Open Paper} action.

The backend separates each logical paper from its observed revisions. Derived sections, chunks, embeddings, summaries, and citations retain revision and transformation identity, while durable idempotent jobs make document processing recoverable. Paper-specific question answering performs exact revision-scoped vector retrieval, applies a fixed lexical reranker, refuses generation below an evidence threshold, and checks generated citation markers against the supplied excerpts. The public source-match state reports lexical support; semantic entailment remains outside this metric. Behavioral adaptation uses versioned events, exposure-qualified negative feedback, temporal decay, per-paper saturation, separate positive and negative centroids, and confidence-weighted recommendation scoring. Recommendation reasons are generated deterministically from the components that participated in ranking. The public citation graph uses resolved paper-to-paper citation relations, bounded traversal, deterministic weighting, clustering, ranking, and a fallback state. The final implementation adds paper-scoped follow-up questions using bounded conversation history and prepares missing citation neighborhoods on demand from cached observations, scholarly metadata providers, or explicit bibliography references.

The Android client is implemented with Kotlin, Jetpack Compose, Room, WorkManager, Retrofit, and a locally packaged D3 graph renderer. ViewModels expose typed live, cached, processing, fallback, empty, and error states. Room stores briefing and paper snapshots together with a duplicate-safe event queue. WorkManager retries interrupted uploads, and server-side event UUIDs prevent repeated delivery from applying the same behavioral signal twice. The graph renderer reports node selection through a restricted JavaScript bridge, while the native layer preserves selection, relation context, navigation, and the action that opens a selected paper. The completed MVP adds a Saved collection, public paper sharing, weekly-briefing notifications, relation and subject filters, and graph recentering through the same client architecture.

The evaluation combines software tests, controlled scenario experiments, live provider runs, and mobile measurements. A coverage-instrumented backend run comprised 533 passing tests, with 89.05\% line and 76.15\% branch coverage; the expanded final regression suite comprised 716 tests and also passed in full. Across eleven applicable controlled behavior traces, the proposed model achieved a mean nDCG@5 of 0.92 and a mean reciprocal rank of 0.94. The recency-only baseline achieved 0.92 and 0.88 respectively. Mechanism comparisons showed the intended effects of exposed negative feedback, exposure gating, per-paper saturation, and duplicate replay. A separate explicit-interest experiment exercised topic addition, replacement, and deletion with eight equal-age candidates. All four directional checks passed; replacement and deletion moved the affected topic group's mean rank by two positions and updated the recommendation reasons. A route-level closed-loop test also changed the top result and its reason after the stored topic was replaced. Aggregate results did not isolate an independent ranking gain for dual-timescale decay or confidence gating.

The live evidence-grounded question evaluation contained fifteen cases over seven papers. Every answered case carried a citation that passed the lexical source check, none of the twelve answerable cases was falsely refused, and two of the three unanswerable cases followed the required refusal-marker protocol. Dense retrieval reached the section designated by the case author in 4 of 12 answerable cases. A model-tier comparison found weaker refusal-marker adherence on the cheaper route. In a separate seven-paper summary comparison, that route satisfied all measured structural criteria and reduced measured cost from \$0.921 to \$0.140. These criteria measure schema completion and lexical source anchoring. Human assessment of explanatory quality remains future evaluation.

The Android resource evaluation retained 80 startup-and-recovery timings and 320 graph timings across eight separately cold-booted Android~14 emulator sessions. Cached cold-launch medians ranged from 0.84 to 1.47~s across the four CPU-and-RAM conditions, while foreground-recovery medians ranged from 66.68 to 93.96~ms. The subsequent graph display in the same WebView had a lower median than the first display in all sixteen resource-by-node cells. First initialization from three fixed arXiv seeds displayed five live papers in every trial, with seed-to-briefing latency ranging from 28.46 to 41.73~s and a median of 36.50~s. Ten sequential event-synchronization batches preserved pending state through database reopening and an HTTP~503 response. The isolated backend accepted all 30 unique event identifiers and identified all 30 subsequent replays as duplicates. The interest editor retained added, edited, and removed topics across navigation, while client data-layer tests verified server normalization, local persistence, briefing refresh order, and failure handling. A complete live path produced a five-paper briefing from a seed paper, opened paper detail, answered a free-form question with sources, rendered a multi-node citation graph, and opened a selected neighbor. The final integrated session additionally saved and shared the selected paper, completed a source-linked follow-up, filtered a graph of 50 papers and 53 directed citations, and opened the corresponding ten-paper weekly briefing from its notification; all 32 connected-device regression tests passed.

Together, these results show that the implemented provenance, evidence, behavior, and mobile-integration mechanisms operate through one continuous literature workflow. The evaluation provides reproducible evidence from software tests, live model runs, controlled behavior scenarios, client measurements, and cross-layer execution. Its scope comprises fifteen question cases, lexical source matching, designed behavior traces, one Android version, and one emulator host. The implemented graph controls are verified by software tests and await participant retesting. Broader studies are required for research productivity, scientific correctness, long-term recommendation quality, and cross-device performance.

\end{digest}
